With this letter, I am submitting my application for the 
% job title
postdoctoral research position on metabolic modeling of the rumen microbiome to guide methane reduction strategies.
I am currently a 
% current position
postdoctoral researcher at the \href{http://www.msysbiology.com/}{Lab of Microbial Systems Biology}, KU Leuven, working on the~\href{https://www.3domics.eu/}{3D'omics} project, that focuses on poultry and swine gut microbiome. 
% role
Specifically, I am involved in the development of microbial network models that incorporate spatial and metabolic data to predict species interactions and evaluate how metabolites influence species distributions.

My research interests focus on unraveling the \textit{Gordian knot} of microbial ecosystems,
the laws that govern their structure and dynamics and the role of evolution in shaping them. 
In the loop of communities shaped by and re-shaping back their habitat, microbial interactions work as the interface where the sum effects of those laws are reflected. 
Therefore, my ambition is to hack microbial interactions in terms of developing a framework for their interpretation and dynamics, and I am confident that metabolic modeling, when integrated with real-world data, can play a pivotal role in advancing this endeavor.

During my PhD, I focused on developing workflows and software tools tailored to real-world microbial communities, particularly for the analysis of high-throughput sequencing data. 
I also served as a lead developer of a \href{https://prego.hcmr.gr/}{data integration platform} designed to extract and consolidate associations among species, environments, and biological processes from both literature and abundance tables from key platforms such as MGnify, and IMG/JGI.

Over the last two years, I have been developing a \href{https://microbetag.readthedocs.io/}{microbial co-occurrence network annotator tool}, and deepened my understanding of pairwise metabolic complementarities, particularly within the context of cross-feeding interactions, through topological analysis of their corresponding metabolic models. 
Further, I have been contributing on the \href{https://geomscale.github.io/}{GeomScale project}
developing an \href{https://github.com/geomScale/dingo}{efficient framework for flux sampling} within the solution space of metabolic models.
To support further our metabolic modeling based efforts, I have also been working on 
the development of a \href{https://mgrowthdb.gbiomed.kuleuven.be/}{microbial growth data repository}; a key resource for constraining models more accurately. 

% why i am interested in the specific position
The H2Rumen project and the associated position strongly align with my research interests and would provide an excellent opportunity to continue my work on microbial interactions through metabolic modeling. 
The methods developed to address the project's challenges would also be highly relevant to my own scientific interests, particularly the interplay between resource hierarchies and species-specific auxotrophies. 
This alignment ensures a strong personal motivation to contribute to the project's success.

\newpage\vspace*{0.3in} 

I am confident that I can apply my current expertise to contribute effectively to the project’s goals, while also acquiring new skills in the process—an aspect I find particularly valuable.
I am enclosing my CV and a letter of recommendation from Prof. Faust. 
Please do not hesitate to contact me for anything that might be needed. 

